Con \(t\) en minutos, el volumen es \(V(t)=50+2t\) litros. El dosificador debe aportar
\[
r(t)=1+0.10t\quad\mathrm{g/min}.
\]
La capacidad del depósito permite llegar hasta \(t=15\) min, mientras que la del dosificador permite seguir la concentración solicitada sólo hasta \(t=10\) min, inclusive. Limita primero el dosificador; entonces hay \(70\,\mathrm L\), a \(0.20\,\mathrm{g/L}\).
